\documentclass[11pt,spanish]{article}
\usepackage[utopia]{mathdesign}
\renewcommand{\familydefault}{\rmdefault}
\usepackage[T1]{fontenc}
\usepackage[utf8]{inputenc}
\usepackage{geometry}
\geometry{verbose,tmargin=2.5cm,bmargin=2.5cm,lmargin=2.5cm,rmargin=2.5cm}
\usepackage{fancyhdr}
\pagestyle{fancy}
\setlength{\parskip}{\bigskipamount}
\setlength{\parindent}{0pt}
\usepackage{array}
\usepackage{longtable}
\usepackage{pmboxdraw}
\usepackage{multirow}
\usepackage{graphicx}

\providecommand{\tabularnewline}{\\}
\usepackage{titling}
\setlength{\droptitle}{-2cm}
\usepackage{multicol}
\usepackage{enumitem}
\usepackage[tikz]{bclogo}

\definecolor{anti-flashwhite}{rgb}{0.95, 0.95, 0.96}
\definecolor{celestialblue}{rgb}{0.29, 0.59, 0.82}
\definecolor{ao(english)}{rgb}{0.0, 0.5, 0.0}

\usepackage{babel}
\addto\shorthandsspanish{\spanishdeactivate{~<>.}}

\usepackage{listings}
\renewcommand{\lstlistingname}{Código}

\begin{document}

\lhead{IP 2026-1}

\rhead{Estructura física de una computadora}
\title{Introducción a la Programación, 2026-1\\
Nota de clase 02: Estructura física de una computadora}
\author{Manuel Soto Romero}
\date{\today \\ Colegio de Ciencia y Tecnología UACM}

\maketitle
\setcounter{section}{2}

En esta nota revisaremos, de manera general, cómo se organiza la computadora
para almacenar y procesar información. Presentaremos definiciones básicas
de hardware, software y programa, introduciremos la arquitectura de von
Neumann y describiremos el ciclo de ejecución de instrucciones con un
ejemplo sencillo.

La programación de computadoras se ha vuelto una labor indispensable
hoy en día. Podemos encontrar programas en cualquier dispositivo,
celulares, tabletas, computadoras personales, televisiones, entre
muchos otros. A lo largo de estas notas estudiaremos los conceptos
básicos que involucran a la programación de computadoras y nos enfocaremos
principalmente en el diseño de algoritmos que puedan ser automatizados.

En esta nota, en particular, discutiremos de forma muy general cómo
es que la computadora almacena información, es decir, estudiaremos
(de forma muy superficial) su arquitectura. Esto con el fin de entender,
a grandes rasgos, lo que sucede cuando un programa se ejecuta. Estudiaremos
con detalle la llamada \emph{Arquitectura de von Neumann} en la cual
se basan las computadoras actuales.

\subsection*{\color{celestialblue}█ \quad{}\color{black}Definiciones básicas}

Comencemos por dar una definición de lo que es una computadora.

\begin{bclogo}[noborder=true, couleurBarre=celestialblue,arrondi =0.1, couleur =anti-flashwhite, logo=\bcplume, marge =10]{\textit{Definición} (\textbf{Computadora})}
\vspace{0.25cm}

\emph{Es una máquina electrónica que trabaja a gran velocidad. Procesa
datos mediante lenguajes entendibles para la misma, almacena y brinda
información.}

\vspace{0.25cm}       
\end{bclogo}

A simple vista, una computadora no es más que un conjunto de componentes
físicos. Llamamos a estos componentes el hardware de la computadora.

\begin{bclogo}[noborder=true, couleurBarre=celestialblue,arrondi =0.1, couleur =anti-flashwhite, logo=\bcplume, marge =10]{\textit{Definición} (\textbf{Hardware})}
\vspace{0.25cm}

\emph{Es la estructura física de una computadora, todo lo tangible,
es decir, los dispositivos materiales que puede tocar la persona usuaria como
son el teclado, el monitor, el gabinete, etcétera.}

\vspace{0.25cm}       
\end{bclogo}

El hardware de una computadora se encarga de procesar datos, almacenar
y brindar información. Sin embargo, la mayor utilidad
de una computadora proviene del hecho de que son entes programables,
o dicho de otro modo, una computadora es útil porque ejecuta software.

\begin{bclogo}[noborder=true, couleurBarre=celestialblue,arrondi =0.1, couleur =anti-flashwhite, logo=\bcplume, marge =10]{\textit{Definición} (\textbf{Software})}
\vspace{0.25cm}

\emph{Es la estructura lógica de una computadora, todo lo intangible,
es decir, son las aplicaciones, sistemas operativos, etcétera que
no son tangibles para la persona usuaria.}

\vspace{0.25cm}       
\end{bclogo}

Esta relación entre hardware y software se logra gracias a la programación
de computadoras. En sus inicios, esta labor consistía en reconfigurar
el hardware, desconectando cables y conectándolos en otros lugares.
Hoy en día la programación de computadoras se lleva a cabo usando
lo que se conoce como \emph{lenguaje de alto nivel}.

\subsection*{\color{celestialblue}█ \quad{}\color{black}Arquitectura de von
Neumann}

Conscientes de que la labor de programación era muy lenta y, por lo
tanto, propensa a errores, J. Presper Eckert, John Mauchly, Herman
Goldstine y John von Neumann discutieron acerca de una forma de proveer
a la computadora del programa a ejecutar de manera más sencilla.
Con lo cual se definió una arquitectura con el siguiente proceso:
\begin{enumerate}
\item Cargar el programa en la memoria de la computadora.
\item La persona usuaria proporciona datos de entrada a la computadora.
\item A partir de la información de entrada, el programa es ejecutado.
\item Una vez ejecutado el programa, se obtienen datos de salida.
\end{enumerate}
Dentro de este proceso, intervienen los siguientes componentes y unidades:
\cite{key-2}
\begin{itemize}
\item \textbf{Memoria RAM}

Aquí son almacenados los datos y las instrucciones durante la ejecución
del programa.
\item \textbf{Unidad Central de Proceso (}\textbf{\emph{CPU, Central Processing
Unit}}\textbf{)}

Controla y coordina la ejecución de las instrucciones de un programa.
Se compone de dos unidades auxiliares: 
\begin{itemize}
\item La \emph{Unidad Aritmético-Lógica }(ALU) se encarga de realizar operaciones
aritméticas (sumas y restas) y lógicas (Álgebra Booleana) sobre los
datos del programa. 
\item La \emph{Unidad de Control} se encarga de leer las instrucciones del
programa, decodificarlas y enviar órdenes a los componentes correspondientes.
Cada instrucción es guardada temporalmente en lo que se conoce como
\emph{Registro de Instrucción} y el \emph{Contador del Programa} que
indica cuál es la siguiente instrucción a ejecutar.
\end{itemize}
\item \textbf{Dispositivos de entrada y salida}

Son responsables de la comunicación con la persona usuaria. Los dispositivos
de entrada permiten a la computadora obtener datos e instrucciones,
por ejemplo, a través del teclado. Los dispositivos de salida, permiten
enviar los resultados de un programa a la persona usuaria, por ejemplo, a través
del monitor.
\item \textbf{Buses}

Permiten conectar todos los componentes. Son vistos como una línea
de transmisión donde pasa la información. Existen, en general tres
tipos de \emph{buses:}
\begin{itemize}
\item \emph{Bus de datos}. Es un canal \emph{bidireccional }que\emph{ }permite
el intercambio de datos entre la CPU y el resto de unidades.
\item \emph{Bus de direcciones}. Es un canal \emph{unidireccional} por el
cual la unidad central de proceso envía las direcciones de memoria
para ubicar información en los dispositivos de memoria.
\item \emph{Bus de control}. Por el bus de control, se transmiten las órdenes
necesarias para activar las operaciones de lectura o escritura.
\end{itemize}
\end{itemize}
\begin{bclogo}[noborder=true, couleurBarre=celestialblue,arrondi =0.1, couleur =anti-flashwhite, logo=\bclampe, marge =10]{\textit{Observación}}
\vspace{0.25cm}

Es importante mencionar que la Unidad Central de Proceso no recibe
los datos de manera directa de la unidad de entrada, ni los envía
directamente a la unidad de salida. Usa la memoria RAM como lugar
de paso. De esta forma:
\begin{itemize}
\item La Unidad Central de Proceso es la única que puede procesar los datos,
esto implica que los datos tienen que llegar de alguna forma ahí para
ser procesados.
\item La Unidad Central de Proceso sólo puede acceder a los datos o instrucciones
almacenados en memoria RAM.
\end{itemize}
\vspace{0.25cm}       
\end{bclogo}

\begin{bclogo}[noborder=true, couleurBarre=celestialblue,arrondi =0.1, couleur =anti-flashwhite, logo=\bclampe, marge =10]{\textit{Observación}}
\vspace{0.25cm}

Un registro es, de forma más detallada, una colección de celdas de
memoria, cada una de las cuales puede almacenar un 0 o un 1 (bit).
El número de celdas corresponde al número de bits que pueden almacenarse
en un registro. De esta forma hay registros de 4, 16 o 32 bits. A
esto se le conoce como \emph{palabra}.

\vspace{0.25cm}       
\end{bclogo}

La Figura 1 muestra un diagrama de la llamada \emph{Arquitectura de
von Neumann}.
\begin{center}
\begin{figure}[h]
\begin{centering}
\includegraphics[scale=0.7]{imagenes/1}
\par\end{centering}
\caption{\emph{Arquitectura de von Neumann\cite{key-1}}}
\end{figure}
\par\end{center}

Finalmente, John von Neumann publicó esta arquitectura en el artículo
\emph{First Draft of a Reporter: on the EDVAC} a partir de donde recibió
este nombre. Esta arquitectura es tomada como base en la mayoría de
dispositivos actuales.

De esta forma, el funcionamiento general de una computadora, se obtiene
mediante un programa (software), o dicho de otro modo, sigue los siguientes
pasos:
\begin{center}
\emph{}%
\begin{tabular}{cc>{\centering}p{5cm}cc}
Datos de entrada & $\Rightarrow$ & Computadora & $\Rightarrow$ & Datos de salida\tabularnewline
 &  & $\Uparrow$ &  & \tabularnewline
 &  & Programa (Conformado por instrucciones) &  & \tabularnewline
\end{tabular}
\par\end{center}

\subsection*{\color{celestialblue}█ \quad{}\color{black}Procesamiento de instrucciones}

Hemos dicho que una computadora toma datos de entrada que son procesados
mediante un programa, generando datos de salida. Revisemos la definición
de estos conceptos con un poco más de detenimiento:

\begin{bclogo}[noborder=true, couleurBarre=celestialblue,arrondi =0.1, couleur =anti-flashwhite, logo=\bcplume, marge =10]{\textit{Definición} (\textbf{Dato})}
\vspace{0.25cm}

\emph{Un dato es un símbolo, valor, texto o fecha que da lugar a la
información. Un dato por sí mismo no nos proporciona contexto alguno; sin
embargo, con el procesamiento adecuado da lugar a obtener información. Se
dice que un dato es la unidad mínima de información.}

\vspace{0.25cm}       
\end{bclogo}

\begin{bclogo}[noborder=true, couleurBarre=celestialblue,arrondi =0.1, couleur =anti-flashwhite, logo=\bccrayon, marge =10]{\textit{Ejemplo}}
\vspace{0.25cm}
\begin{itemize}
\item 'F'
\item ``Chilaquiles''
\item 2001
\item 4/02/2004
\end{itemize}
\vspace{0.25cm}       
\end{bclogo}

\begin{bclogo}[noborder=true, couleurBarre=celestialblue,arrondi =0.1, couleur =anti-flashwhite, logo=\bcplume, marge =10]{\textit{Definición} (\textbf{Información})}
\vspace{0.25cm}

\emph{Es un conjunto de datos procesados en un determinado contexto.}

\vspace{0.25cm}       
\end{bclogo}

\begin{bclogo}[noborder=true, couleurBarre=celestialblue,arrondi =0.1, couleur =anti-flashwhite, logo=\bccrayon, marge =10]{\textit{Ejemplo}}
\vspace{0.25cm}
\begin{itemize}
\item 'F' Representa el género femenino. 
\item ``Chilaquiles'' Es la comida favorita de la abuelita de Batman. 
\item 2001 Es el año de fundación de la UACM. 
\item 4/02/2004 Es la fecha de fundación de \textsc{Facebook}.
\end{itemize}
\vspace{0.25cm}       
\end{bclogo}

\subsubsection*{Programas}

Debido a que las computadoras únicamente entienden lenguaje de máquina
(binario), fue necesario crear nuevos lenguajes que fueran más entendibles
para las personas y que funcionaran como un intermediario entre la
persona programadora y la computadora. Este tipo de lenguajes (de programación)
son clasificados como \emph{lenguajes de alto nivel }y \emph{lenguajes
de bajo nivel}. Mientras más cercano a la arquitectura de la computadora
(hardware) más bajo el nivel.
\begin{itemize}
\item De esta forma, lenguajes de programación como \textsc{C}, \textsc{Java}
o \textsc{Python} que son más cercanos al lenguaje de los humanos,
son de alto nivel.
\item Los lenguajes de bajo nivel son llamados \emph{ensambladores} y cada
procesador provee su propio lenguaje.
\end{itemize}
Con estos conceptos definidos, veamos ahora, qué es un programa.

\begin{bclogo}[noborder=true, couleurBarre=celestialblue,arrondi =0.1, couleur =anti-flashwhite, logo=\bcplume, marge =10]{\textit{Definición} (\textbf{Programa})}
\vspace{0.25cm}

\emph{Es un conjunto de instrucciones que describen las órdenes dadas
a una computadora, éstas pueden ser escritas en un lenguaje de programación,
de alto o bajo nivel o bien en lenguaje máquina.}

\vspace{0.25cm}       
\end{bclogo}
\begin{itemize}
\item Los programas escritos en lenguaje máquina son ejecutados directamente
por la computadora.
\item Los programas escritos en lenguajes de programación deben ser traducidos
al lenguaje de la máquina para que ésta pueda ejecutarlos.
\end{itemize}
Al programa que se requiere traducir se le llama \emph{programa fuente}\footnote{También llamado código fuente.}.
Para realizar la traducción, se puede emplear un programa llamado
\emph{traductor}. En general existen dos tipos de traductores: \emph{compiladores}
e \emph{intérpretes}. El programa final generado después de la traducción
es llamado \emph{programa objeto}\footnote{Código objeto}.

\begin{bclogo}[noborder=true, couleurBarre=celestialblue,arrondi =0.1, couleur =anti-flashwhite, logo=\bcplume, marge =10]{\textit{Definición} (\textbf{Compilador})}
\vspace{0.25cm}

\emph{Es un traductor de lenguajes de programación de alto o bajo
nivel que toma todo un programa fuente y lo traduce pasando por distintas
fases en un programa ensamblador o programa objeto.}

\vspace{0.25cm}       
\end{bclogo}

\begin{bclogo}[noborder=true, couleurBarre=celestialblue,arrondi =0.1, couleur =anti-flashwhite, logo=\bcplume, marge =10]{\textit{Definición} (\textbf{Intérprete})}
\vspace{0.25cm}

\emph{Es un traductor de lenguajes de programación de alto o bajo
nivel que traduce línea por línea las instrucciones de un programa
fuente, es decir, traduce cada línea a código objeto y la ejecuta.}

\vspace{0.25cm}       
\end{bclogo}

\subsubsection*{Representación de la información}

Los datos e instrucciones son almacenados en la memoria en un lenguaje
entendible para la misma, llamado \emph{lenguaje de máquina}. Este
lenguaje se compone de cadenas de ceros y unos y es llamado \emph{binario}.
Una computadora entiende este lenguaje, pues únicamente trabaja con
dos estados: \emph{activo o desactivado, abierto o cerrado, pasa corriente
o no pasa corriente}. 

\bigskip{}

\textbf{Unidades de medida}

Tanto la memoria como las unidades de almacenamiento utilizan unidades
para ser medidas, entre las cuales se encuentran.
\begin{itemize}
\item \textbf{Bit.} Es la unidad más pequeña de información en el sistema
binario, es decir, es o bien 0 o 1.
\item \textbf{Byte.} Es conocida como la unidad de almacenamiento, se presenta
en bloques de 8 bits. Con esto se pueden representar hasta 256 valores
distintos. 
\item \textbf{Kilobyte (KB).} Corresponde a 1,024 bytes.
\item \textbf{Megabyte (MB).} Corresponde a 1,024 kilobytes, es decir, 1,048,576
bytes.
\item \textbf{Gigabyte (GB). }Corresponde a 1,024 megabytes, es decir, 1,073,741,824
bytes.
\item \textbf{Terabyte (TB). }Corresponde a 1,024 gigabytes, es decir, 1,099,511,627,776
bytes.
\end{itemize}

\subsubsection*{Ejecución de instrucciones}

Ahora que tenemos bien definido cómo se representa la información
en la computadora, y cómo es traducida a lenguaje de máquina, veamos
cómo es que la computadora procesa cada cadena de bits.

La mayoría de los lenguajes de programación se traducen directamente
al lenguaje que provee el procesador, llamado ensamblador. Por lo
que ejemplificaremos el proceso de ejecución, mediante programas escritos
en este tipo de lenguajes. Por ejemplo, supongamos un lenguaje ensamblador
con las siguientes operaciones: \cite{key-2}
\begin{center}
\begin{longtable}[c]{lc>{\raggedright}p{10cm}}
\textbf{\emph{Instrucción}} & \textbf{\emph{Código}} & \textbf{\emph{Descripción}}\tabularnewline
\hline 
\endhead
\texttt{\textbf{Carga <DIR>}} & 20 & Transfiere el dato almacenado en la posición de memoria \texttt{\textbf{<DIR>}}
al acumulador de la ALU.\tabularnewline
\hline 
\texttt{\textbf{Guarda <DIR>}} & 02 & Guarda el contenido del acumulador en la posición de memoria \texttt{\textbf{<DIR>}}.\tabularnewline
\hline 
\texttt{\textbf{Divide <DIR>}} & 38 & Divide el contenido del acumulador entre el dato almacenado en la
posición de memoria \texttt{\textbf{<DIR>}}. El resultado queda almacenado
en el acumulador.\tabularnewline
\hline 
\texttt{\textbf{Multiplica <DIR>}} & 36 & Multiplica el contenido del acumulador entre el dato almacenado en
la posición de memoria \texttt{\textbf{<DIR>}}. El resultado queda
almacenado en el acumulador.\tabularnewline
\hline 
\texttt{\textbf{Resta <DIR>}} & 33 & Resta el contenido del acumulador entre el dato almacenado en la posición
de memoria \texttt{\textbf{<DIR>}}. El resultado queda almacenado
en el acumulador.\tabularnewline
\hline 
\texttt{\textbf{Suma <DIR>}} & 30 & Suma el contenido del acumulador entre el dato almacenado en la posición
de memoria \texttt{\textbf{<DIR>}}. El resultado queda almacenado
en el acumulador.\tabularnewline
\hline 
\texttt{\textbf{Alto}} & 70 & Indica el fin del programa\tabularnewline
\end{longtable}
\par\end{center}

\begin{bclogo}[noborder=true, couleurBarre=celestialblue,arrondi =0.1, couleur =anti-flashwhite, logo=\bccrayon, marge =10]{\textit{Ejemplo}}
\vspace{0.25cm}

Supongamos que tenemos un programa escrito en un lenguaje de alto
nivel, que calcula el promedio de tres números. Algo similar al siguiente
código:

\bigskip{}

\begin{lstlisting}[basicstyle={\ttfamily},tabsize=4,mathescape=true]
Imprimir "Ingresa el primer numero"
Leer numero1
Imprimir "Ingresa el segundo numero"
Leer numero2
Imprimir "Ingresa el tercer numero"
Leer numero3
promedio $\leftarrow$ (numero1 + numero2 + numero3) / 3
Imprimir "Promedio: " + promedio
\end{lstlisting}
Supongamos que la persona usuaria ingresa los números, \texttt{\textbf{10}},
\texttt{\textbf{4}} y \texttt{\textbf{4}}, por lo tanto, las variables
\texttt{numero1}, \texttt{numero2} y \texttt{numero3}, tienen esos
valores, que podemos almacenar en los registros \texttt{\textbf{00}},
\texttt{\textbf{01}} y \texttt{\textbf{02}} de la memoria. En el registro
\texttt{\textbf{03}}, agrego un número 3 que me permitirá hacer la
división. Con lo cual el código ensamblador correspondiente quedaría
como sigue:

\bigskip{}

\begin{lstlisting}[basicstyle={\ttfamily},tabsize=4,mathescape=true]
Carga  00
Suma   01
Guarda 04
Carga  02
Suma   04
Guarda 04
Carga  04
Divide 03
Guarda 04
Alto
\end{lstlisting}
Una vez realizado esto, necesitamos buscar el código asociado a cada
operación y representarlo en código máquina, es decir, en binario. Esto
se vería así:

\bigskip{}

\begin{lstlisting}[basicstyle={\ttfamily},tabsize=4,mathescape=true]
00010100 00000000
00011110 00000001
00000010 00000100
00010100 00000010
00011110 00000100
00000010 00000100
00010100 00000100
00100110 00000011
00010100 00000100
01000110 00000000
\end{lstlisting}
Los pasos que sigue el procesador para ejecutar instrucciones se conocen
como \emph{ciclo de ejecución }y se listan a continuación\cite{key-2}:
\begin{enumerate}
\item La unidad de control obtiene de la memoria RAM la instrucción indicada
en el contador de programa (PC) y la almacena en el registro de instrucción
(IR). Actualiza el contenido del contador de programa con la dirección
de memoria donde se encuentra almacenada la siguiente instrucción
del programa.
\item La unidad de control decodifica la instrucción del registro de instrucción,
detecta las unidades que deben interactuar para ejecutarla y les envía
señales de control para indicarles la acción que deben de realizar.
Si la acción requiere un dato que esté en memoria, se hará la lectura
correspondiente, utilizando para ello los buses apropiados. 
\item Si las instrucciones contienen operaciones lógicas o aritméticas en
los datos, este trabajo es asignado a la ALU, la cual realiza la operación
con el dato indicado en la instrucción y el dato almacenado en el
acumulador. El resultado es almacenado en el acumulador de la ALU, perdiéndose
con ello el dato anterior.
\end{enumerate}
De esta forma, la memoria final de nuestro programa, luce de la siguiente
manera. Suponiendo que las instrucciones se almacenan a partir de
la dirección en memoria 10.
\begin{center}
\begin{tabular}{|>{\centering}p{1.5cm}|>{\centering}p{1.5cm}|>{\centering}p{1.5cm}|>{\centering}p{1.5cm}|c|>{\centering}p{1.5cm}|>{\centering}p{1.5cm}|c}
\cline{1-4} \cline{2-4} \cline{3-4} \cline{4-4} \cline{6-7} \cline{7-7} 
00000000

00001010 & 0000000

00000100 & 0000000

00000100 & 0000000

0000011 & ... & 00010100

00000000 & 00011110

00000001 & ...\tabularnewline
\cline{1-4} \cline{2-4} \cline{3-4} \cline{4-4} \cline{6-7} \cline{7-7} 
\multicolumn{1}{>{\centering}p{1.5cm}}{00} & \multicolumn{1}{>{\centering}p{1.5cm}}{01} & \multicolumn{1}{>{\centering}p{1.5cm}}{02} & \multicolumn{1}{>{\centering}p{1.5cm}}{03} & \multicolumn{1}{c}{} & \multicolumn{1}{>{\centering}p{1.5cm}}{10} & \multicolumn{1}{>{\centering}p{1.5cm}}{11} & \tabularnewline
\end{tabular}
\par\end{center}

Podemos resumir la ejecución de cada instrucción en la siguiente tabla.
\cite{key-2}
\begin{center}
\begin{longtable}[c]{|l|c|c|c|c|c|c|c|c|c|c|}
\hline 
\multirow{2}{*}{Instrucción} & \multicolumn{2}{c|}{ALU} & \multicolumn{3}{c|}{Unidad de control} & \multicolumn{5}{c|}{Memoria}\tabularnewline
\cline{2-6} \cline{3-6} \cline{4-6} \cline{5-6} \cline{6-6} 
 & Operación & Acumulador & \multicolumn{2}{c|}{IR} & PC & \multicolumn{5}{c|}{}\tabularnewline
\hline 
\multirow{2}{*}{} & \multirow{2}{*}{} & \multirow{2}{*}{} & \multirow{2}{*}{} & \multirow{2}{*}{} & \multirow{2}{*}{10} & 10 & 4 & 4 & 3 & \tabularnewline
\cline{7-11} \cline{8-11} \cline{9-11} \cline{10-11} \cline{11-11} 
 &  &  &  &  &  & \texttt{\textbf{00}} & \texttt{\textbf{01}} & \texttt{\textbf{02}} & \texttt{\textbf{03}} & \texttt{\textbf{04}}\tabularnewline
\hline 
\multirow{2}{*}{\texttt{\textbf{Carga 00}}} & \multirow{2}{*}{} & \multirow{2}{*}{10} & \multirow{2}{*}{20} & \multirow{2}{*}{00} & \multirow{2}{*}{11} & 10 & 4 & 4 & 3 & \tabularnewline
\cline{7-11} \cline{8-11} \cline{9-11} \cline{10-11} \cline{11-11} 
 &  &  &  &  &  & \texttt{\textbf{00}} & \texttt{\textbf{01}} & \texttt{\textbf{02}} & \texttt{\textbf{03}} & \texttt{\textbf{04}}\tabularnewline
\hline 
\multirow{2}{*}{\texttt{\textbf{Suma 01}}} & \multirow{2}{*}{\texttt{\textbf{10+4}}} & \multirow{2}{*}{14} & \multirow{2}{*}{30} & \multirow{2}{*}{01} & \multirow{2}{*}{12} & 10 & 4 & 4 & 3 & \tabularnewline
\cline{7-11} \cline{8-11} \cline{9-11} \cline{10-11} \cline{11-11} 
 &  &  &  &  &  & \texttt{\textbf{00}} & \texttt{\textbf{01}} & \texttt{\textbf{02}} & \texttt{\textbf{03}} & \texttt{\textbf{04}}\tabularnewline
\hline 
\multirow{2}{*}{\texttt{\textbf{Guarda 04}}} & \multirow{2}{*}{} & \multirow{2}{*}{14} & \multirow{2}{*}{02} & \multirow{2}{*}{04} & \multirow{2}{*}{13} & 10 & 4 & 4 & 3 & 14\tabularnewline
\cline{7-11} \cline{8-11} \cline{9-11} \cline{10-11} \cline{11-11} 
 &  &  &  &  &  & \texttt{\textbf{00}} & \texttt{\textbf{01}} & \texttt{\textbf{02}} & \texttt{\textbf{03}} & \texttt{\textbf{04}}\tabularnewline
\hline 
\multirow{2}{*}{\texttt{\textbf{Carga 02}}} & \multirow{2}{*}{} & \multirow{2}{*}{4} & \multirow{2}{*}{20} & \multirow{2}{*}{02} & \multirow{2}{*}{14} & 10 & 4 & 4 & 3 & 14\tabularnewline
\cline{7-11} \cline{8-11} \cline{9-11} \cline{10-11} \cline{11-11} 
 &  &  &  &  &  & \texttt{\textbf{00}} & \texttt{\textbf{01}} & \texttt{\textbf{02}} & \texttt{\textbf{03}} & \texttt{\textbf{04}}\tabularnewline
\hline 
\multirow{2}{*}{\texttt{\textbf{Suma 04}}} & \multirow{2}{*}{4+14} & \multirow{2}{*}{18} & \multirow{2}{*}{30} & \multirow{2}{*}{04} & \multirow{2}{*}{15} & 10 & 4 & 4 & 3 & 14\tabularnewline
\cline{7-11} \cline{8-11} \cline{9-11} \cline{10-11} \cline{11-11} 
 &  &  &  &  &  & \texttt{\textbf{00}} & \texttt{\textbf{01}} & \texttt{\textbf{02}} & \texttt{\textbf{03}} & \texttt{\textbf{04}}\tabularnewline
\hline 
\multirow{2}{*}{\texttt{\textbf{Guarda 04}}} & \multirow{2}{*}{} & \multirow{2}{*}{18} & \multirow{2}{*}{02} & \multirow{2}{*}{04} & \multirow{2}{*}{16} & 10 & 4 & 4 & 3 & 18\tabularnewline
\cline{7-11} \cline{8-11} \cline{9-11} \cline{10-11} \cline{11-11} 
 &  &  &  &  &  & \texttt{\textbf{00}} & \texttt{\textbf{01}} & \texttt{\textbf{02}} & \texttt{\textbf{03}} & \texttt{\textbf{04}}\tabularnewline
\hline 
\multirow{2}{*}{\texttt{\textbf{Carga 04}}} & \multirow{2}{*}{} & \multirow{2}{*}{18} & \multirow{2}{*}{20} & \multirow{2}{*}{04} & \multirow{2}{*}{17} & 10 & 4 & 4 & 3 & 18\tabularnewline
\cline{7-11} \cline{8-11} \cline{9-11} \cline{10-11} \cline{11-11} 
 &  &  &  &  &  & \texttt{\textbf{00}} & \texttt{\textbf{01}} & \texttt{\textbf{02}} & \texttt{\textbf{03}} & \texttt{\textbf{04}}\tabularnewline
\hline 
\multirow{2}{*}{\texttt{\textbf{Divide 03}}} & \multirow{2}{*}{\texttt{\textbf{18/3}}} & \multirow{2}{*}{6} & \multirow{2}{*}{38} & \multirow{2}{*}{03} & \multirow{2}{*}{18} & 10 & 4 & 4 & 3 & 18\tabularnewline
\cline{7-11} \cline{8-11} \cline{9-11} \cline{10-11} \cline{11-11} 
 &  &  &  &  &  & \texttt{\textbf{00}} & \texttt{\textbf{01}} & \texttt{\textbf{02}} & \texttt{\textbf{03}} & \texttt{\textbf{04}}\tabularnewline
\hline 
\multirow{2}{*}{\texttt{\textbf{Guarda 04}}} & \multirow{2}{*}{} & \multirow{2}{*}{6} & \multirow{2}{*}{20} & \multirow{2}{*}{04} & \multirow{2}{*}{19} & 10 & 4 & 4 & 3 & 6\tabularnewline
\cline{7-11} \cline{8-11} \cline{9-11} \cline{10-11} \cline{11-11} 
 &  &  &  &  &  & \texttt{\textbf{00}} & \texttt{\textbf{01}} & \texttt{\textbf{02}} & \texttt{\textbf{03}} & \texttt{\textbf{04}}\tabularnewline
\hline 
\multirow{2}{*}{\texttt{\textbf{Alto}}} & \multirow{2}{*}{} & \multirow{2}{*}{} & \multirow{2}{*}{70} & \multirow{2}{*}{} & \multirow{2}{*}{-} & 10 & 4 & 4 & 3 & 6\tabularnewline
\cline{7-11} \cline{8-11} \cline{9-11} \cline{10-11} \cline{11-11} 
 &  &  &  &  &  & \texttt{\textbf{00}} & \texttt{\textbf{01}} & \texttt{\textbf{02}} & \texttt{\textbf{03}} & \texttt{\textbf{04}}\tabularnewline
\hline 
\end{longtable}
\par\end{center}

\vspace{0.25cm}       
\end{bclogo}

\subsection*{\color{celestialblue}█ \quad{}\color{black}Conclusiones}

Comprender la arquitectura básica de una computadora ayuda a conectar los
conceptos de programación con el comportamiento real del hardware. Al
identificar el papel de la memoria, la CPU y los buses, así como el ciclo
de ejecución de instrucciones, podemos razonar mejor sobre cómo se
transforman los datos en resultados y por qué el diseño de algoritmos y
programas es tan importante.
\begin{thebibliography}{1}
\bibitem{key-1}José de J. Galaviz, \emph{Notas de Clase de Organización
y Arquitectura de Computadoras}, Facultad de Ciencias UNAM, Revisión
2015-2.

\bibitem{key-2}Alejandra Andrade, Diana Cruz, et. al. \emph{Programación
estructurada y modular en el lenguaje C}, UACM, Primera Edición, 2010. 

\bibitem{key-3}Juan J. Alvarez, S. Alejandra Andrade, et. al. \emph{Introducción
a la Programación}, UACM, Primera edición 2007.
\end{thebibliography}

\end{document}
